% from ega4-21c-fr.tex
\oldpage[IV-4]{309}
Supposons qu'il existe dans $X_y$ une composante irréductible $Z$ dont le point
générique $z$ n'est pas isolé dans $Z$, et telle que $\operatorname{codim}(Z,X)\geq2$.
La question étant locale sur $X$ et $Y$ puisque $z$ est non isolé dans $Z$, on peut, en
remplaçant $X$ par un voisinage ouvert de $z$ dans $X$, supposer $X$ et $Y$ affines,
$X_y=Z$ irréductible, non réduit à un point et telle que
$\operatorname{codim}(X_y,X)\geq2$. L'image $g(X-X_y)=U$ est un ouvert de $Y'$
isomorphe à $X-X_y$ par hypothèse; montrons qu'en remplaçant au besoin $X$ par un
voisinage ouvert de $z$ dans $X$, on peut supposer que $U$ ne contient aucune des
composantes irréductibles de $Y'$ dont l'adhérence contient $y$.

En effet, on peut tout d'abord supposer que tous les points maximaux de $X$ sont des
générisations de $z$, l'ensemble de ces points étant fini; donc l'ensemble des points
maximaux de $U$ est l'ensemble des images $y_i$ par $g$ des points maximaux $x_i$ de
$X$ (aucun point maximal de $X$ ne pouvant par hypothèse être contenu dans $X_y$,
puisque $\dim(\mathcal O_{X,z})\geq2$). Posons $X_i=\overline{\{x_i\}}$. Par
hypothèse, on a $z\in X_i$ et puisque $z$ n'est pas isolé dans $X_y$, il existe dans
$X_y$ un point $x\ne z$. Comme $X_i\ne Z$, il existe dans $X_i$ un point $t_i$ tel
que si l'on pose $T_i=\overline{\{t_i\}}$, on ait $\dim(\mathcal O_{T_i,x})=1$
(10.5.9); cela entraîne par hypothèse $t_i\notin Z$, donc $t_i$ n'est pas une
générisation de $z$. Remplaçant $X$ par $X'=X-\bigcup_iT_i$, on voit que l'image par
$g$ de $X'-X_y$ ne contient pas les images $g(t_i)$, donc ne contient aucune des
composantes irréductibles de $Y'$ dont l'adhérence contient $y$.

Cela étant, comme $X$ et $Y$ sont affines, le morphisme $g:X\to Y$ est affine, donc il
en est de même de la restriction $g':X'\to Y'$; puisque $g'$ est une immersion ouverte,
l'immersion canonique $U\to Y'$ est affine. Posons
$Y_1=\operatorname{Spec}(\mathcal O_{Y,y})$, $Y'_1=Y_1-\{y\}=Y'\cap Y_1$,
$U_1=U\cap Y_1$. Ce qui précède prouve que l'immersion canonique $U_1\to Y'_1$ est
affine, donc $U_1$ est un ouvert affine dans $Y_1$ en vertu de (W), en d'autres termes
l'immersion canonique $U_1\to Y_1$ est affine. Mais puisque cette immersion est de
présentation finie (1.6.2), il résulte de (8.10.5, (viii)) appliqué suivant la méthode de
(8.1.2, a)) qu'en restreignant au besoin $Y$ à un voisinage ouvert affine de $y$, on peut
supposer que l'immersion $U\to Y$ est affine. On en conclurait que l'ouvert $X-X_y$ de
$X$, isomorphe à $U$, serait affine, ce qui contredirait (21.12.7); la proposition est
ainsi démontrée.
\end{proof}

\begin{corollary}[21.12.11]
\label{IV.21.12.11}
Soient $X$, $Y$ deux préschémas localement noethériens, $X$ étant supposé
irréductible; soit $g:X\to Y$ un morphisme localement de type fini, et soit $V$ le plus
grand ouvert de $X$ tel que la restriction $g|V:V\to Y$ soit un isomorphisme local.
Supposons que pour tout point $y\in Y$, l'anneau local $\mathcal O_{Y,y}$ vérifie la
condition (W) (21.12.8). Alors toute composante irréductible de $T=X-V$ est, ou bien
de codimension $\leq1$, ou bien telle que son point générique $z$ soit isolé dans
$g^{-1}(g(z))$.
\end{corollary}

\begin{proof}
Posons $g(z)=y$. La question ne dépendant que de la fibre $g^{-1}(y)$ et de l'anneau
$\mathcal O_{Y,y}$, on peut par changement de base se borner au cas où
$Y=\operatorname{Spec}(\mathcal O_{Y,y})$ et où $X$ est affine
$((\mathbf I, 3.6.5)$ et $(\mathbf I, 4.5.5))$; remplaçant éventuellement $X$ par un
voisinage ouvert de $z$ dans $X$, on peut supposer que $T$ est irréductible; en outre on
peut se borner au cas où $V$ est non vide. Le morphisme $g$ peut donc être supposé
séparé et $y$ fermé dans $Y$; puisque la restriction $g|V:V\to Y$ est un isomorphisme
local et que l'ouvert non vide $V$ de $X$ est irréductible, il résulte de
$(\mathbf I, 8.2.8)$ que $g|V:V\to Y$ est une immersion ouverte. La restriction
$g|(V\cap X_y):V\cap X_y\to\{y\}=\operatorname{Spec}(k(y))$ est donc aussi une
immersion ouverte, ce qui montre que $V\cap X_y$ est, soit vide, soit réduit à un point
$x$ rationnel
